% ===================== Codificacion e idioma =====================
\usepackage[utf8]{inputenc}
\usepackage[T1]{fontenc}
\usepackage[spanish,es-noquoting,es-nolists,es-noshorthands]{babel}
\usepackage{lmodern}
\usepackage{microtype}

% ===================== Matematicas =====================
\usepackage{amsmath,amssymb,amsthm,mathtools}
\usepackage{bm}

% ===================== Formato =====================
\usepackage[letterpaper,margin=2.6cm,bottom=3cm,headheight=16pt]{geometry}
\usepackage{setspace}
\usepackage{graphicx}
\usepackage{booktabs}
\usepackage{longtable}
\usepackage{array}
\usepackage{multirow}
\usepackage{float}
\usepackage[font=small,labelfont=bf,width=.92\textwidth]{caption}
\usepackage{subcaption}
\usepackage{enumitem}
\usepackage{xcolor}
\usepackage{fancyhdr}
\usepackage[ruled,vlined,linesnumbered,spanish]{algorithm2e}
\usepackage{tikz}
\usetikzlibrary{arrows.meta,positioning,calc,decorations.pathreplacing,shapes.geometric}
\usepackage[round,authoryear]{natbib}
\usepackage{xurl}
\usepackage[hidelinks,breaklinks]{hyperref}
\usepackage[nameinlink,spanish]{cleveref}


% ===================== Listados de codigo fuente =====================
\usepackage{listings}

% listings no maneja UTF-8 por si mismo con pdflatex: hay que declarar
% explicitamente cada caracter acentuado. Sin esto, los comentarios en
% espanol del codigo fuente salen truncados o con caracteres perdidos.
\lstset{
  inputencoding=utf8,
  extendedchars=true,
  literate=%
    {á}{{\'a}}1 {é}{{\'e}}1 {í}{{\'i}}1 {ó}{{\'o}}1 {ú}{{\'u}}1
    {Á}{{\'A}}1 {É}{{\'E}}1 {Í}{{\'I}}1 {Ó}{{\'O}}1 {Ú}{{\'U}}1
    {ñ}{{\~n}}1 {Ñ}{{\~N}}1 {ü}{{\"u}}1 {Ü}{{\"U}}1
    {¿}{{?`}}1 {¡}{{!`}}1 {°}{{\textdegree}}1 {º}{{\textordmasculine}}1
    {–}{{--}}1 {—}{{---}}1 {“}{{``}}1 {”}{{''}}1 {‘}{{`}}1 {’}{{'}}1
    {…}{{\ldots}}1 {≤}{{$\leq$}}1 {≥}{{$\geq$}}1 {×}{{$\times$}}1
    {Δ}{{$\Delta$}}1 {∂}{{$\partial$}}1 {α}{{$\alpha$}}1 {λ}{{$\lambda$}}1
    {μ}{{$\mu$}}1 {φ}{{$\varphi$}}1 {ρ}{{$\rho$}}1 {σ}{{$\sigma$}}1
    {τ}{{$\tau$}}1 {θ}{{$\theta$}}1 {γ}{{$\gamma$}}1 {δ}{{$\delta$}}1
    {ε}{{$\varepsilon$}}1 {∞}{{$\infty$}}1 {≠}{{$\neq$}}1 {→}{{$\to$}}1
    {∈}{{$\in$}}1 {·}{{\textperiodcentered}}1
}

\definecolor{codfondo}{HTML}{FAFAFA}
\definecolor{codcomentario}{HTML}{5E7A87}
\definecolor{codclave}{HTML}{0057A3}
\definecolor{codcadena}{HTML}{9B5000}
\definecolor{codnumero}{HTML}{9AA5AD}
\definecolor{codmarco}{HTML}{DEE3E6}

\lstdefinestyle{python}{
  language=Python,
  basicstyle=\ttfamily\scriptsize,
  keywordstyle=\color{codclave}\bfseries,
  commentstyle=\color{codcomentario}\itshape,
  stringstyle=\color{codcadena},
  numberstyle=\ttfamily\tiny\color{codnumero},
  numbers=left, numbersep=7pt,
  backgroundcolor=\color{codfondo},
  frame=single, framerule=.4pt, rulecolor=\color{codmarco},
  breaklines=true, breakatwhitespace=false,
  postbreak=\mbox{\textcolor{codnumero}{$\hookrightarrow$}\space},
  showstringspaces=false,
  tabsize=4,
  captionpos=b,
  xleftmargin=1.6em, xrightmargin=.2em,
  aboveskip=1.1em, belowskip=1.1em,
  morekeywords={self,True,False,None,dataclass,field,replace},
}
\lstset{style=python}
\renewcommand{\lstlistingname}{Listado}
\renewcommand{\lstlistlistingname}{Índice de listados}

% ===================== Colores =====================
\definecolor{azultesis}{HTML}{0072B2}
\definecolor{naranja}{HTML}{E69F00}
\definecolor{verdetesis}{HTML}{009E73}
\definecolor{rojotesis}{HTML}{D55E00}
\definecolor{grisclaro}{HTML}{F2F2F2}
\definecolor{grismedio}{HTML}{7F7F7F}

\hypersetup{colorlinks=true, linkcolor=azultesis, citecolor=verdetesis,
            urlcolor=azultesis, linktoc=all}

% ===================== Encabezados =====================
\pagestyle{fancy}
\fancyhf{}
\renewcommand{\headrulewidth}{0.3pt}
\fancyhead[L]{\small\nouppercase{\leftmark}}
\fancyhead[R]{\small\thepage}
\fancypagestyle{plain}{\fancyhf{}\renewcommand{\headrulewidth}{0pt}\fancyfoot[C]{\small\thepage}}

% ===================== Entornos de teorema =====================
\theoremstyle{plain}
\newtheorem{theorem}{Teorema}[chapter]
\newtheorem{proposition}[theorem]{Proposición}
\newtheorem{lemma}[theorem]{Lema}
\newtheorem{corollary}[theorem]{Corolario}
\theoremstyle{definition}
\newtheorem{definition}[theorem]{Definición}
\newtheorem{example}[theorem]{Ejemplo}
\newtheorem{assumption}[theorem]{Hipótesis}
\theoremstyle{remark}
\newtheorem{remark}[theorem]{Observación}
\newtheorem{criticism}[theorem]{Discusión crítica}
\crefname{criticism}{discusión crítica}{discusiones críticas}
\Crefname{criticism}{Discusión crítica}{Discusiones críticas}
\crefname{remark}{observación}{observaciones}
\Crefname{remark}{Observación}{Observaciones}
% Ecuaciones: solo el número entre paréntesis, p. ej. «de (3.11)».
\crefformat{equation}{(#2#1#3)}
\Crefformat{equation}{(#2#1#3)}
\crefrangeformat{equation}{(#3#1#4) a~(#5#2#6)}
\crefmultiformat{equation}{(#2#1#3)}{ y~(#2#1#3)}{, (#2#1#3)}{ y~(#2#1#3)}
\Crefmultiformat{equation}{(#2#1#3)}{ y~(#2#1#3)}{, (#2#1#3)}{ y~(#2#1#3)}

% Caja para los resultados originales del trabajo
\usepackage[most]{tcolorbox}
\newtcolorbox{aporte}[1][]{
  colback=grisclaro, colframe=azultesis, boxrule=.6pt, arc=1.6pt,
  left=6pt,right=6pt,top=5pt,bottom=5pt, #1}

% ===================== Notacion =====================
\newcommand{\R}{\mathbb{R}}
\newcommand{\N}{\mathbb{N}}
\newcommand{\Prob}{\mathbb{P}}
\newcommand{\Esp}{\mathbb{E}}
\newcommand{\ind}{\mathbf{1}}
\newcommand{\dB}{\partial_B}
\newcommand{\dC}{\partial}
\newcommand{\KA}{\mathit{KA}}
\newcommand{\VP}{\mathit{VP}}
\newcommand{\CONAB}{\mathit{CONAB}}
\newcommand{\SMLV}{\mathit{SMLV}}
\newcommand{\AUF}{\mathit{AUF}}
\newcommand{\RAB}{\mathit{RAB}}
\newcommand{\CAI}{\mathit{CAI}}
\newcommand{\DURA}{\mathit{DURA}}
\newcommand{\npx}[2]{{}_{#1}p_{#2}}
\newcommand{\nqx}[2]{{}_{#1}q_{#2}}
\newcommand{\med}{\operatorname{med}}
\newcommand{\sgn}{\operatorname{sgn}}
\newcommand{\dom}{\operatorname{dom}}
\newcommand{\Lip}{\operatorname{Lip}}
\DeclareMathOperator*{\argmin}{arg\,min}

% Espaciado de tablas
\renewcommand{\arraystretch}{1.15}
\setlength{\tabcolsep}{5pt}

% Titulos de capitulo compactos
\usepackage{titlesec}
\titleformat{\chapter}[display]
  {\normalfont\huge\bfseries\color{azultesis}}
  {\chaptertitlename\ \thechapter}{14pt}{\Huge}
\titlespacing*{\chapter}{0pt}{-18pt}{28pt}

% Enlace al repositorio de la tesis (codigo, tablas completas y aplicacion).
\newcommand{\repositorio}{\url{https://github.com/elagosd/tesis-rentas-vitalicias}}
